% Generated draft: requires mathematical and source review.
\documentclass[11pt]{article}
\usepackage[margin=25mm]{geometry}
\usepackage{fontspec}
\setmainfont{Latin Modern Roman}
\usepackage{amsmath,amssymb,mathtools,unicode-math}
\setmathfont{Latin Modern Math}
\usepackage{enumitem,needspace,xurl,hyperref}
\setlength{\parindent}{0pt}
\setlength{\parskip}{5pt}
\setlength{\emergencystretch}{3em}
\newcommand{\N}{\mathbb{N}}
\newcommand{\Z}{\mathbb{Z}}
\newcommand{\Q}{\mathbb{Q}}
\newcommand{\R}{\mathbb{R}}
\newcommand{\C}{\mathbb{C}}
\newcommand{\F}{\mathbb{F}}
\newcommand{\A}{\mathbb{A}}
\newcommand{\PP}{\mathbb{P}}
\newcommand{\E}{\mathbb{E}}
\newcommand{\eps}{\varepsilon}
\newcommand{\dd}{\,\mathrm{d}}
\begin{document}
\section*{UnsolvedMath ID: 800918}
\textbf{Source problem code:} catalogue entry
\subsection{Definitions and mathematical statement}
The supplied question concerns a space denoted by \(M^{K3}_{2n}\), whose moduli convention and ground field are not specified in the record. A usual interpretation is a moduli space or stack of polarized K3 surfaces \((S,L)\) with \(L^2=2n\), for a positive integer \(n\). This interpretation is conditional: the source excerpt does not specify whether the polarization is primitive, whether it is ample or merely quasi-polarizing, or whether a stack or coarse moduli space is intended.

For a geometric K3 surface \(S\), its Picard lattice is \(\operatorname{Pic}(S)\) equipped with the intersection pairing, and its Picard rank is
\[
\rho(S)=\operatorname{rank}\operatorname{Pic}(S).
\]
Noether--Lefschetz loci parameterize surfaces whose Picard lattices have higher rank than the generic Picard lattice in the intended moduli problem. In the usual polarized setting with generic Picard rank one, these are loci with \(\rho(S)\geq 2\), with individual loci specified by additional lattice data.

After fixing the moduli convention, an ambient intersection or cohomology ring \(A^*(M^{K3}_{2n})\), its coefficient ring, and a collection \(\mathcal N\) of Noether--Lefschetz loci admitting classes in that ring, define the candidate
\[
R_{\mathrm{NL}}^{A}(M^{K3}_{2n})
=
\bigl\langle [Z]\;\bigm|\;Z\in\mathcal N\bigr\rangle_{A^*(M^{K3}_{2n})},
\]
where the right-hand side denotes the unital subring generated by the indicated classes.

\textbf{Question (Pandharipande, 2010, Q6).} What is the analogue of the tautological ring for \(M^{K3}_{2n}\)? The supplied source proposes the subring generated by classes of Noether--Lefschetz loci as a natural candidate. It does not specify an independent definition or characterization of the desired tautological ring, an ambient ring, coefficients, or the precise collection of loci. Thus the question is a request to identify and characterize an appropriate ring; the excerpt alone does not determine a precise equality conjecture.
\subsection{Short English statement}
Identify an appropriate tautological ring for the moduli of K3 surfaces denoted by \(M^{K3}_{2n}\). Pandharipande suggests the ring generated by Noether--Lefschetz classes, which describe loci where the Picard rank increases. The precise moduli convention, ambient ring, coefficients, and intended characterization require clarification.
\subsection{Sources}
\begin{itemize}
\item Local UnsolvedMath registry, permanent ID \(800918\), exact source problem code \texttt{catalogue entry}; title: AGNES Open Problems in Algebraic Geometry --- Q6. The supplied record attributes the question to Pandharipande, proposed in 2010, and labels it open. This status is supplied registry data, not an independently verified literature assessment.
\item Supplied source list: Pandharipande, \emph{Open Problems in Enumerative Algebraic Geometry (AGNES)} (2010), source item Q6 (\#5). Supplied source URL: \url{https://arxiv.org/abs/1004.3259}. The background records access on 2026-10-16 and extraction by \texttt{arxiv\_latex}; these are supplied provenance statements, not a claim of independent source access or verification.
\item Supplied upstream revision: \texttt{d371ee917ac9647fbb06cd6fc529deec6f3dda82}.
\end{itemize}
\end{document}
